{\small
\begin{quote}
	\textit{
		[C]uidado con el matem\'atico y los otros que hacen falsas profec\'{\i}as.
		El peligro ya existe de que los matem\'aticos hayan hecho un pacto con el
		Diablo para oscurecer el esp\'{\i}ritu y confinar al Hombre dentro de las
		ataduras del Infierno.
	}
	\hfill ---San Agustín
\end{quote}
}
\paragraph{Preguntas para investigar}
\begin{enumerate}[label=P\arabic*]
	
	\item\label{item:factorizacion}
		?`Qu\'e se entiende por la \emph{descomposici\'on canónica}
		de un entero en primos positivos?
	\item\label{item:criba}
		Usa la criba de Erat\'ostenes para hacer una lista de todos
		los primos menores que $200$.
	\seti
\end{enumerate}
%

\paragraph{Una pregunta para explorar}
\begin{enumerate}[label=P\arabic*]
	\conti
	\item\label{item:gaussianos}
		Considera los n\'umeros complejos $a + bi$ con $a$ y $b$ en
		$\Enteros$ (donde $\raizcuarta = \sqrt{-1}$).
		Estos n\'umeros forman bajo la adici\'on y multiplicaci\'on
		un sistema matem\'atico no muy diferente de $\Enteros$.
		Este sistema com\'unmente se denota por $\EnterosGauss$
		¿Cu\'ales de las propiedades en P7--P12 de la Lista~\#1
		y en P4--P7 de la Lista~\#2 se cumplen en $\EnterosGauss$?
	\seti
\end{enumerate}
%

\paragraph{Probar o refutar, y arreglar si es posible}
Haz una lista de propiedades de $\Enteros$ que creas que da
una buena descripci\'on de este sistema matemático
¿Tu lista es suficiente para deducir todo lo que has probado hasta
ahora sobre $\Enteros$?
\begin{enumerate}[label=P\arabic*]
	\conti
	\item\label{item:division:orden}
		Si $a$ y $b$ son enteros positivos y $a \mid b$ entonces $a \leq b$.
	\item\label{item:division:cancelacion}
		Si $a$, $b$  $c$ son enteros y $ab \mid ac$, entonces $b \mid c$.
	\item\label{item:multiplicacion:negativo}
		Si $a > b$ y $c < 0$ entonces $ac < bc$.
		Verdadero en $\Enteros$.
	\item\label{item:division:suma}
		Si $a \mid b$ y $a \mid (b + c)$, entonces $a \mid c$.
	\item\label{item:division:algoritmo}
		Dados cualquier par de n\'umeros naturales $a, b \in \mathbb{N}$,
		existen enteros $q$ y $r$ tales que $a = bq + r$ y $0 \leq r < b$.
	\item\label{item:division:alternativa}
		Si $a \mid bc$ y $a \nmid b$ entonces $a \mid c$.
	\item\label{item:primos:sucesion}
		Si $p_1 = 2$, $p_2 = 3$, $p_3 = 5$, \ldots, $p_n$ son los primeros $n$
		n\'umeros primos, entonces $P_n = p_1 p_2 p_3 \cdots p_n + 1$
		es primo para todo $n$.
	\seti
\end{enumerate}
%

\paragraph{Problemas num\'ericos}
Algo para reflexionar.
\begin{enumerate}[label=P\arabic*]
	\conti
	\item\label{item:modulares:unidades}
		Escribe todos los elementos de $\Unidadesmod[22]$,
		de $\Unidadesmod[23]$, de $\Unidadesmod[25]$,
		de $\Unidadesmod[105]$
		¿Cu\'ales tienen ra\'{\i}ces primitivas?
		¿Alguna conjetura?
	\item\label{item:modulares:logaritmos}
		Construye una tabla de logaritmos para $\Unidadesmod[19]$.
		Encuentra todos los cuadrados perfectos en $\Unidadesmod[19]$.
	\item\label{item:fracciones:procedimiento}
		Describe (y justifica, si es posible) el proceso representado
		por la siguiente tabla para calcular las convergentes sucesivas
		en P9 de la fracción continua en P10 de la Lista~\#2.
		\begin{center}
			\begin{tabular}{cc|c|c|c|c|c|}
				& & 2 & 1 & 1 & 1 & 3  \\
				\hline
				0 & 1 & 2 & 3 & 5 & 8 & 29 \\ \hline 
				1 & 0 & 1 & 1 & 2 & 3 & 11 \\
				\hline	
			\end{tabular}
			%
		\end{center}
		?`Notaste algo interesante sobre los determinantes de $2 \times 2$
		\begin{displaymath}
			\begin{vmatrix} 1 & 2 \\ 0 & 1 \end{vmatrix}, \quad
			\begin{vmatrix} 2 & 3 \\ 1 & 1 \end{vmatrix}, \quad
			\begin{vmatrix} 3 & 5 \\ 1 & 2 \end{vmatrix}, \quad
			\begin{vmatrix} 5 & 8 \\ 2 & 3 \end{vmatrix}, \quad
			\begin{vmatrix} 8 & 29 \\ 3 & 11 \end{vmatrix}
				\text{ ?}
		\end{displaymath}
		%
	\item\label{item:fracciones:convergencia}
		Procede con la fracci\'on $\frac{7469}{2463}$ como en P8, P9 y P10 de la Lista~\#2.
		Grafica los valores de todas las convergentes en la recta real y observa
		cuidadosamente las relaciones entre las convergentes y la fracci\'on misma,
		as\'{\i} como las relaciones entre las convergentes entre s\'{\i}.
		Al aproximar la fracci\'on mediante las convergentes de su expansi\'on
		en fracci\'on continua,
		¿el precio a pagar por la precisi\'on en cada caso es inesperadamente alto? ¿Bajo?
	\item\label{item:ecuaciones:lineales}
		Usa tus resultados de los dos \'ultimos problemas para encontrar una soluci\'on
		entera para cada una de las ecuaciones diofánticas lineales:
		$29x + 11y = 1$ y $7469x + 2463y = 1$.
	\seti
\end{enumerate}
%

\paragraph{Ingenio}
\begin{enumerate}[label=P\arabic*]
	\conti
	\item\label{item:reticulo}
		Considera un sistema de coordenadas cartesianas rectangulares en un plano.
		A los puntos con coordenadas enteras los llamamos \emph{puntos reticulares}.
		Demuestra que ning\'un tri\'angulo cuyos v\'ertices sean todos puntos reticulares
		puede ser un tri\'angulo equil\'atero.
	\seti
\end{enumerate}
%

